\section{Stochastic calculus for asset prices}

Three ideas carry continuous-time pricing: Brownian motion and its quadratic variation, It\^o's lemma for functions of a diffusing price, and the representation of a price as a discounted expectation under a changed measure. The option models that use them are in \hyperref[mod:black-scholes-merton]{the Black--Scholes--Merton model}; numerical solution of the resulting equations is in \hyperref[meth:finite-difference-option-pricing]{finite-difference pricing}.

\begin{qfdefinition}[Brownian motion and quadratic variation: independent Gaussian increments, quadratic variation equal to elapsed time, and why $(dW)^2=dt$ produces It\^o's correction]\label{def:brownian-motion-quadratic-variation}
\qffield{Statement} A Wiener process (Brownian motion) $z$ has two properties: over a short interval $\Delta t$ its change is $\Delta z=\epsilon\sqrt{\Delta t}$ with $\epsilon$ standard normal, and changes over non-overlapping intervals are independent. So $z(T)-z(0)$ is normal with mean zero and variance $T$, and uncertainty grows with the square root of the horizon. Because $\sqrt{\Delta t}$ is much larger than $\Delta t$ for small steps, paths are jagged; their expected length over any interval is infinite \citep[chap.~14, sec.~14.2, pp.~318--319]{s7-hulloptions}. Wilmott builds the same process as the limit of a fair coin-tossing game with $n$ bets of size $\sqrt{t/n}$, and lists its properties: finiteness, continuity, the Markov and martingale properties, normal increments, and quadratic variation $\sum_{j=1}^n(X(t_j)-X(t_{j-1}))^2\to t$ \citep[chap.~4, secs.~4.5--4.6, pp.~73--75]{wide-wilmott-pwoqf}.

\emph{Why it matters.} Each squared increment has mean $\Delta t$ and variance $2\Delta t^2$, so the sum of $n$ of them over $[0,t]$ has mean $t$ and variance $2t^2/n$. In the limit the randomness disappears and $\int_0^t(dX)^2=t$ in the mean-square sense, which is the rule $dX^2=dt$ \citep[secs.~4.9 and~4.12.1, pp.~77--78 and~82--83]{wide-wilmott-pwoqf}. Hull's argument is the same: $\epsilon^2\Delta t$ has expected value $\Delta t$ and variance of order $\Delta t^2$, too small to be stochastic in the limit (appendix to chap.~14, pp.~336--337). A second-order term that ordinary calculus discards therefore survives as a drift.

\qffield{Notation} $W_t$ (Hull's $z$, Wilmott's $X$) standard Brownian motion with $W_0=0$; $[W]_T=\lim\sum(W_{t_j}-W_{t_{j-1}})^2$ its quadratic variation over a partition of $[0,T]$ whose mesh goes to zero. Stochastic integrals $\int_0^tf\,dW=\lim\sum f(t_{j-1})(W_{t_j}-W_{t_{j-1}})$ evaluate the integrand at the left point of each interval, so the integrand does not anticipate the increment that multiplies it (Wilmott, sec.~4.7, pp.~75--76).

\qffield{Conditions} Model mathematics, not an empirical claim about prices: the properties hold for the model process. Quadratic variation is a pathwise limit; on a finite grid it is only an estimate with sampling error $\sqrt{2/n}\,T$. Fractional Brownian motion, used in rough-volatility models, gives up independent increments, and with them this scaling (Hull, sec.~14.8, pp.~329--330).

\qffield{Practice} Example, synthetic: one simulated path on $[0,1]$ with 10{,}000 steps has quadratic variation 1.0052, within 0.4 standard deviations (0.0141) of 1. Refining the same path to 20{,}000 steps by Brownian-bridge midpoints gives 0.9996. The first variation $\sum|\Delta W|$ is 79.7, close to $\sqrt{2n/\pi}=79.8$, so it grows without bound as the grid is refined. The left-point sum $\sum W_{t_{j-1}}\Delta W_j$ equals $\tfrac12(W_T^2-\sum\Delta W^2)$ exactly, $-0.2109$ here, near the limit $\tfrac12(W_T^2-T)=-0.2083$. The midpoint sum gives $\tfrac12W_T^2$, with no $-T/2$ term: the choice of evaluation point changes the integral.

\qffield{Pitfalls} Treating $dW$ as small in the way $dt$ is, and dropping $(dW)^2$. Reading the scaling of Brownian increments, variance in proportion to elapsed time and standard deviation in proportion to its square root, as a property of observed returns without testing it (serial correlation breaks it). Using a midpoint or forward-looking integrand in a trading simulation: it uses information about the increment that a position taken at the start of the interval could not have.

\qffield{See also} \hyperref[res:ito-lemma-geometric-brownian-motion]{It\^o's lemma}; \hyperref[meth:monte-carlo-simulation]{Monte Carlo simulation}; \hyperref[def:financial-distributions]{distributions} (the lognormal model).

\qffield{Further reading} \nolinkurl{crypto-shreve} (Volume~II, rigorous construction of the It\^o integral).

\end{qfdefinition}

\begin{qfresult}[It\^o's lemma and geometric Brownian motion: the one-dimensional lemma, the lognormal solution, and mean against median of the terminal price]\label{res:ito-lemma-geometric-brownian-motion}
\qffield{Statement} If $dx=a(x,t)\,dt+b(x,t)\,dz$, a function $G(x,t)$ follows
\[ dG=\Big(\frac{\partial G}{\partial x}a+\frac{\partial G}{\partial t}+\frac12\frac{\partial^2G}{\partial x^2}b^2\Big)dt+\frac{\partial G}{\partial x}b\,dz, \]
driven by the same Wiener process \citep[chap.~14, sec.~14.6, eq.~(14.12), p.~327]{s7-hulloptions}. Wilmott derives it from a Taylor series in which the squared increments sum to $dt$, and gives the rule of thumb: expand to second order and set $dX^2=dt$ \citep[chap.~4, secs.~4.10 and~4.12, pp.~78--83]{wide-wilmott-pwoqf}. For $dS=\mu S\,dt+\sigma S\,dz$ and $G=\ln S$, $\partial G/\partial S=1/S$ and $\partial^2G/\partial S^2=-1/S^2$. So $d\ln S=(\mu-\sigma^2/2)\,dt+\sigma\,dz$ and $\ln S_T$ is normal with mean $\ln S_0+(\mu-\sigma^2/2)T$ and variance $\sigma^2T$ (Hull, eqs.~(14.17)--(14.19), pp.~328--329). Integrating gives $S_T=S_0\exp((\mu-\tfrac12\sigma^2)T+\sigma W_T)$ (Wilmott, sec.~4.14.2, p.~85).

\emph{Mean and median.} The expected price is $E(S_T)=S_0e^{\mu T}$ (Hull, chap.~15, eq.~(15.4), p.~340), while the median is $S_0e^{(\mu-\sigma^2/2)T}$, because the median of a lognormal variable is the exponential of the mean of its log. The expected continuously compounded return is $\mu-\sigma^2/2$, not $\mu$: $\mu\Delta t$ is close to the arithmetic mean of short-period returns, while compounding over the whole period gives a geometric average (sec.~15.3, pp.~341--342).

\qffield{Notation} $\mu$ expected return (drift) and $\sigma$ volatility, both per year and constant; $\epsilon$ standard normal. Interpretation: for $G=S^2$, It\^o gives $d(S^2)=(2\mu+\sigma^2)S^2\,dt+2\sigma S^2\,dz$, so $E[S_T^2]=S_0^2e^{(2\mu+\sigma^2)T}$; the ordinary chain rule would omit $\sigma^2$.

\qffield{Conditions} One state variable and a smooth $G$ (twice differentiable in $x$, once in $t$). Constant $\mu$ and $\sigma$ make GBM's solution explicit. With several variables the lemma adds cross terms $\rho b_1b_2\,\partial^2V/\partial S_1\partial S_2$ \citep[sec.~4.13, eq.~(4.9), p.~84]{wide-wilmott-pwoqf}. GBM is a model: it has no jumps, no volatility clustering and thin tails relative to daily equity returns.

\qffield{Practice} Example, synthetic: $S_0=100$, $\mu=8$ percent, $\sigma=20$ percent, $T=1$. The mean is 108.33 and the median 106.18. From 200{,}000 exact draws the sample mean is 108.36 (standard error 0.049) and the sample median 106.21 (standard error 0.060), both within one standard error. The simulated $E[S_T^2]$ is 12{,}220 (standard error 11.4) against It\^o's 12{,}214; the chain-rule value 11{,}735 is more than 40 standard errors away. An Euler discretisation of $dS$ on the same Brownian paths has a root-mean-square error to the exact solution of 0.44 at 50 steps and 0.043 at 5{,}000: discretisation error is separate from sampling error.

\qffield{Pitfalls} Calling $\mu$ the expected log return, or compounding $\mu$ to get the median. Using the ordinary chain rule on a function of a diffusing price. Simulating GBM by Euler steps when the exact lognormal draw is available, and then reading the discretisation bias as sampling noise.

\qffield{See also} \hyperref[def:brownian-motion-quadratic-variation]{quadratic variation}; \hyperref[res:risk-neutral-pricing-feynman-kac]{risk-neutral pricing}; \hyperref[def:financial-distributions]{distributions}; \hyperref[res:kelly-growth-optimal-fraction]{Kelly sizing} (log growth).

\end{qfresult}

\begin{qfresult}[Risk-neutral pricing, Girsanov's change of drift and Feynman--Kac: a price as a discounted expectation and its link to the pricing equation]\label{res:risk-neutral-pricing-feynman-kac}
\qffield{Statement} \emph{Market price of risk.} Two derivatives on one underlying variable with $df_i/f_i=m_i\,dt+s_i\,dz$ can be combined into an instantaneously riskless portfolio, which must earn $r$. Hence $(m-r)/s=\lambda$ is the same for all of them; $\lambda$ is the market price of risk. Changing $\lambda$ changes expected growth rates but not volatilities: this is Girsanov's theorem, and choosing $\lambda$ defines the probability measure. With $\lambda=0$ every price grows at $r$, the traditional risk-neutral world \citep[chap.~28, sec.~28.1, eqs.~(28.6)--(28.10), pp.~671--674]{s7-hulloptions}. With the money-market account as numeraire, $f/g$ is a martingale and $f_0=\hat E[e^{-\int_0^Tr\,dt}f_T]$, which reduces to $e^{-rT}\hat E(f_T)$ for a constant rate (sec.~28.4, eqs.~(28.17)--(28.19), p.~677). The BSM equation contains no variable affected by risk preferences, which is why prices can be computed as if investors were risk neutral (chap.~15, sec.~15.7, p.~351).

\emph{Feynman--Kac link.} For $dS=\mu S\,dt+\sigma S\,dX$, $V=e^{-r(T-t)}E[F(S_T)]$ satisfies $V_t+\tfrac12\sigma^2S^2V_{SS}+\mu SV_S-rV=0$, from the backward Kolmogorov equation. Replacing $\mu$ by $r$ gives the Black--Scholes equation, so the option value is the present value of the expected payoff under the risk-neutral walk $dS=rS\,dt+\sigma S\,dX$ \citep[chap.~10, sec.~10.11, eqs.~(10.6)--(10.8), pp.~179--180]{wide-wilmott-pwoqf}.

\qffield{Notation} $P$ real-world measure, $Q$ risk-neutral measure; $W$ a $P$-Brownian motion. Interpretation, a direct calculation for one asset with yield $q$: set $\theta=(\mu-r+q)/\sigma$ and $Z_T=\exp(-\theta W_T-\tfrac12\theta^2T)$. Then $E_P[Z_T]=1$, $\tilde W_t=W_t+\theta t$ is a $Q$-Brownian motion, $dS=(r-q)S\,dt+\sigma S\,d\tilde W$, and $E_Q[X]=E_P[Z_TX]$ for any payoff $X$. Reweighting real-world outcomes by $Z_T$ is the same as simulating with drift $r-q$.

\qffield{Conditions} Frictionless, continuous trading in the underlying and a riskless asset, constant $r$, $q$ and $\sigma$ (or known functions of time) and no arbitrage. Under these conditions the hedge replicates the claim and the measure is unique. With jumps or stochastic volatility the market is incomplete, many measures fit the traded prices, and the expectation depends on which one is chosen (\hyperref[mod:merton-jump-diffusion]{jump diffusion}). The real-world drift $\mu$ still matters for the distribution of hedging profit and loss (\hyperref[res:delta-hedging-volatility-mismatch]{hedging with the wrong volatility}).

\qffield{Practice} Example, synthetic: the call with $S_0=100$, $K=105$, $r=3$ percent, $q=0$, $\sigma=20$ percent and $T=1$ has BSM value 7.1281. Drawing 200{,}000 terminal prices under $\mu=8$ percent and weighting each discounted payoff by $Z_T$ ($\theta=0.25$) gives 7.098 (standard error 0.021), 1.4 standard errors from the formula; the weights average 1.0003. Without the weights, discounting the real-world expected payoff at $r$ gives 9.91, which is not a price. By finite differences, the discounted real-world expectation satisfies the equation with $\mu SV_S$ (residual $2\times10^{-7}$) and fails the Black--Scholes equation (residual $-2.58$); with drift $r$ the expectation is the BSM formula.

\qffield{Pitfalls} Reading the risk-neutral probability as a forecast: the risk-neutral probability of exercise is $N(d_2)$, while the real-world one uses $\mu$ in place of $r$ (Wilmott, sec.~10.12, p.~181). Discounting a real-world expected payoff at the risk-free rate. Applying the formula in an incomplete market without stating which measure was chosen.

\qffield{See also} \hyperref[mod:black-scholes-merton]{BSM}; \hyperref[meth:one-period-binomial-valuation]{one-period binomial valuation} (risk-neutral probability on a tree); \hyperref[meth:monte-carlo-simulation]{Monte Carlo simulation}; \hyperref[meth:finite-difference-option-pricing]{finite differences}.

\qffield{Further reading} \nolinkurl{crypto-shreve} (Volume~II: Girsanov's theorem and Feynman--Kac stated with their technical conditions).

\end{qfresult}
